% Part: first-order-logic
% Chapter: beyond
% Section: intuitionistic-logic

\documentclass[../../../include/open-logic-section]{subfiles}

\begin{document}

\olfileid{fol}{byd}{il}

\olsection{وجدانیت پسند منطق}

دوجے درجے تے اعلیٰ درجے دی منطق دے
برعکس، وجدانیت پسند پہلے درجے دی منطق
کلاسیکی روپ نوں محدود کرکے بندی اے؛
ایہ ہور ”تعمیری“ قسم دے استدلال دا
نمونہ بنان لئی اے۔ ہیٹھلی مثالاں
ایدے پیچھے دیاں کجھ وجوہاں واضح
کر سکدیاں نیں۔

فرض کرو اک دن کوئی بندہ تہاڈے کول
آ کے اعلان کرے کہ اوہنے اک قدرتی
عدد~$x$ معلوم کر لیا اے جیہدی
خاصیت ایہ اے: جے $x$ اولی عدد
ہوئے تاں ریمان دا مفروضہ سچ اے،
تے جے $x$ مرکب عدد ہوئے تاں
ریمان دا مفروضہ جھوٹا اے۔
کمال دی خبر!{} ریمان دا مفروضہ
سچ اے یا نہیں، ایہ ریاضی دے
وڈے کھلے سوالاں وچوں اک اے؛
ایتھے لگدا اے کہ سوال نوں
حسابی جانچ تک گھٹا دتا گیا
اے، یعنی ایہ معلوم کرنا رہ
گیا اے کہ اک خاص عدد اولی
اے یا نہیں۔

$x$ دی جادوئی قدر کیہڑی اے؟
اوہ ایویں بیان کردا اے:
$x$ اوہ قدرتی عدد اے
جیہڑا، جے ریمان دا
مفروضہ سچ ہوئے، $7$
دے برابر اے، تے نہیں
تاں $9$ دے برابر۔

تسی غصے وچ اپنے پیسے
واپس منگدے او۔ کلاسیکی
نقطۂ نظر توں اوپرلی
وضاحت واقعی $x$ دی
اک منفرد قدر متعین
کردی اے؛ پر تہاڈے
لئی مطلوب $x$ دی
اوہ قدر اے جیہڑی
\emph{صاف صاف} دتی
گئی ہوئے۔

ہور، شاید گھٹ بناوٹی،
مثال ویکھو۔ اسیں
جاندے آں کہ غیر
ناطق عدد نوں ناطق
قوت تک چڑھا کے
ناطق نتیجہ مل
سکدا اے؛ مثال
لئی
$\sqrt{2}^2 = 2$۔
پر ایہ گل اینی
صاف نہیں کہ غیر
ناطق عدد نوں
\emph{غیر ناطق}
قوت تک چڑھا کے
ناطق نتیجہ مل
سکدا اے یا نہیں۔
ہیٹھلا قضیہ ہاں
وچ جواب دیندا اے:

\begin{thm}
ایسے غیر ناطق
اعداد $a$ تے
$b$ موجود نیں
جیہناں لئی
$a^b$ ناطق
اے۔
\end{thm}

\begin{proof}
$\sqrt{2}^{\sqrt{2}}$
لو۔ جے ایہ
ناطق ہوئے، تاں
کم ہو گیا:
$a = b = \sqrt{2}$
رکھ لو۔ ورنہ
ایہ غیر ناطق
اے۔ ایس صورت
وچ
\[
(\sqrt{2}^{\sqrt{2}})^{\sqrt{2}} = \sqrt{2}^{\sqrt{2} \cdot
  \sqrt{2}} = \sqrt{2}^2 = 2,
\]
جیہڑا یقیناً
ناطق اے۔
اس لئی ایس
صورت وچ
$a$ نوں
$\sqrt{2}^{\sqrt{2}}$
تے $b$ نوں~$\sqrt 2$
رکھ لو۔
\end{proof}

کی ایہ درست
ثبوت اے؟
زیادہ تر ریاضی
دان سمجھدے نیں
کہ اے۔ پر
فیر وی ایتھے
اک گل کجھ
تسلی نہیں دیندی
اے: اسیں
اک خاصیت
رکھن والے
حقیقی اعداد
دے جوڑے
دی موجودگی
ثابت کر
دتی، پر
ایہ نہیں
دس سکدے
کہ اوہ
\emph{کیہڑا}
جوڑا اے۔
اوہی نتیجہ
ایس طرح وی
ثابت ہو سکدا
اے کہ ثبوت
وچ جوڑا
$a$, $b$
\emph{صاف
دتا} جاوے:
$a = \sqrt{3}$
تے
$b = \log_3 4$
لو۔
% PNB-SEM-033: explicit pair لئی دوواں گواہاں دی irrationality دی source وچ مختصر چھڈی جانچ واضح کیتی۔
ایہ دوویں عدد غیر ناطق نیں۔ تین دے مربع جذر نوں
گھٹ توں گھٹ صورت دی صحیح عددی نسبت منّن نال
نسبت دے دوناں عدد تین نال ونڈن جوگ نکلدے نیں،
جیہڑا ایس گھٹ توں گھٹ صورت دے خلاف اے۔ دوجے
عدد نوں مثبت صحیح اعداد دی نسبت منّن نال تین دی
اک مثبت صحیح عددی گھات، چار دی اک مثبت صحیح عددی
گھات دے برابر ہونی پئے گی؛ پہلی تین نال ونڈن جوگ
اے تے دوجی نہیں، ایس لئی ایہ وی ممکن نہیں۔ فیر
\[
a^b = \sqrt{3}^{\log_3 4} = 3^{1/2 \cdot \log_3 4} = (3^{\log_3
  4})^{1/2} = 4^{1/2}= 2,
\]
کیوں جے
$3^{\log_3 x} = x$۔

وجدانیت پسند
منطق ایہو
جہے استدلال
دا نمونہ
بنان لئی
اے جیہدے
وچ پہلے
ثبوت والی
چال دی
اجازت نہیں۔
ایہ ثابت
کرنا کہ
کوئی $x$
موجود اے
جو~$!A(x)$
پورا کردا
اے، ایس
دا مطلب
اے کہ
تسی اک
خاص~$x$
تے ثبوت
دیو کہ
اوہ $!A$
پورا کردا
اے، جداں
دوجے ثبوت
وچ۔ ایہ
ثابت کرن
لئی کہ
$!A$ یا
$!B$ درست
اے، انہاں
وچوں اک
دا ثبوت
دینا پیندا
اے۔

رسمی طور
اُتے وجدانیت
پسند پہلے
درجے دی
منطق، پہلے
درجے دے
!!a{derivation}
نظام نوں
خاص طریقے
نال محدود
کرن نال
ملدی اے۔
ایس طرح
دوجے تے
اعلیٰ درجے
دی منطق
دے وجدانیت
پسند روپ
وی نیں۔
ریاضی دے
حساب نال
ایہ بس
رسمی استنباطی
نظام نیں؛
پر جداں
پہلاں دسیا،
انہاں دا
مقصد ریاضیاتی
استدلال دی
اک قسم
دا نمونہ
بنانا اے۔
کوئی ایہ
سمجھ سکدا
اے کہ ایہ
اوہ استدلال
اے جیہڑا
ریاضی بارے
کسی خاص
فلسفیانہ
موقف (مثلاً
برووَر دی
وجدانیت)
نال جائز
ٹھہردا اے؛
یا ایہ
اک زیادہ
”ٹھوس“
تے اطمینان
بخش ریاضیاتی
استدلال اے
(بشپ دے
تعمیریت
والے انداز
وچ)؛
تے ایس
اُتے بحث
وی ہو
سکدی اے
کہ رسمی
بیان غیر
رسمی محرک
نوں صحیح
پکڑدا اے
یا نہیں۔
پر ساڈا
فلسفیانہ
موقف کوئی
وی ہوئے،
اسیں وجدانیت
پسند منطق
دا رسمی
طور اُتے
پیش کیتی
منطق ونجوں
مطالعہ کر
سکدے آں؛
تے کئی
ریاضیاتی
منطق دان،
اپنیاں
وجوہاں نال،
ایہ مطالعہ
دلچسپ سمجھدے
نیں۔

وجدانیت پسند
رابطیاں دی
اک غیر رسمی
تعمیری تشریح
اے، جیہنوں
عموماً BHK
تشریح آکھدے
نیں؛ ایہ
نام برووَر،
ہیٹنگ تے
کولموگوروف
دے ناماں
توں بنیا
اے۔ تشریح
ایویں اے:
$!A
\land !B$
دا ثبوت،
$!A$ دے
ثبوت تے
$!B$ دے
ثبوت دی
جوڑی اے؛
$!A \lor !B$
دا ثبوت،
$!A$ یا
$!B$ وچوں
اک دا ثبوت
اے، تے
صاف خبر
ہوئے کہ
کیہڑا؛
$!A \lif !B$
دا ثبوت
اک طریقہ
اے جو
$!A$ دے
ثبوت نوں
~$!B$ دے
ثبوت وچ
بدلدا اے؛
$\lforall[x][!A(x)]$
دا ثبوت
اک طریقہ
اے جو
$!A(x)$ دا
ثبوت دیندا
اے،~$x$
دی ہر قدر
لئی؛
تے
$\lexists[x][!A(x)]$
دا ثبوت
~$x$ دی
اک قدر
تے ایس
قدر دے
~$!A$ پورا
کرن دا
ثبوت اے۔
ایہ تشریح
حسابی اصطلاحاں
وچ ”کری--ہاورڈ
تماثل“ یا
”!!{formula}s
بطور انواع
دا اصول“
وی کہلا
سکدی اے:
!!a{formula}
نوں اک
خاص قسم
دے ڈیٹا
ٹائپ دی
تعیین
سمجھو، تے
ثبوتاں نوں
انہاں ڈیٹا
ٹائپاں دیاں
حسابی شےواں،
جیہڑیاں
متعلقہ
!!{formula}
دی سچائی
سامنے لے
آندیاں
نیں۔

وجدانیت پسند
منطق نوں
اکثر کلاسیکی
منطق ”وچوں“
تیسرے امکان
دے اخراج دا
قانون کڈ
کے باقی
رہن والی
منطق سمجھیا
جاندا اے۔
ہیٹھلا قضیہ
ایہ گل
ہور عین
کردا اے۔
\begin{thm}
وجدانیت پسند
منطق وچ
ہیٹھلے مسلّمی
خاکے اک
دوجے دے
برابر نیں:
\begin{enumerate}
\item $(\lnot !A \lif \lfalse) \lif !A$.
\item $!A \lor \lnot !A$
\item $\lnot \lnot !A \lif !A$
\end{enumerate}
\end{thm}
ایہ برابری سارے خاکیاں نوں اضافی مسلّمے منّن بارے
اے؛ ایہ نہیں کہ دکھائے ہوئے فارمولے وجدانیت پسند
منطق وچ پہلے توں ثابت نیں۔ سارے خاکے دا فرض
کرنا تے اک دتے ہوئے فارمولے لئی اوہدی صرف اک
مخصوص صورت دا فرض کرنا وکھ وکھ گلّاں نیں۔
اک خاکے دی
مخصوص صورت
باقی دوواں
وچوں کسے
اک توں
کڈنا وجدانیت
پسند منطق
دی چنگی
مشق اے۔

وجدانیت پسند
قضیاتی منطق
دے پہلے
استنباطی نظام،
برووَر دی
وجدانیت
دی رسمی
صورتاں ونجوں،
کولموگوروف،
گلیوینکو
تے ہیٹنگ
نے اک دوجے
توں آزاد
پیش کیتے۔
وجدانیت پسند
پہلے درجے
دی منطق
(تے وجدانیت
پسند ریاضی
دے کجھ حصیاں)
دی پہلی
رسمی صورت
ہیٹنگ نے
دتی۔
اگرچہ کئی
کلاسیکی
طور اُتے
درست خاکے
وجدانیت پسند
طور اُتے
درست نہیں،
پر کئی
ہوندے نیں۔

\emph{دوہری نفی دا ترجمہ}
کلاسیکی تے
وجدانیت پسند
منطق دے
اک اہم رشتے
نوں بیان
کردا اے۔
ایہ ہیٹھاں
استقرائی
طور اُتے
متعین اے
(~$!A^N$
نوں کلاسیکی
!!{formula}~$!A$
دا ”وجدانیت
پسند“ ترجمہ
سمجھو):
\begin{align*}
!A^N & \ident \lnot\lnot !A \quad \text{جوہری !!{formula}s $!A$ لئی} \\
(!A \land !B)^N & \ident (!A^N \land !B^N) \\
(!A \lor !B)^N & \ident  \lnot\lnot (!A^N \lor !B^N) \\
(!A \lif !B)^N & \ident (!A^N \lif !B^N) \\
(\lforall[x][!A])^N & \ident \lforall[x][!A^N] \\
(\lexists[x][!A])^N & \ident \lnot\lnot\lexists[x][!A^N]
\end{align*}
کولموگوروف تے
گلیوینکو نے
قضیاتی منطق
لئی ایس ترجمے
دے روپ دتے
سن؛ محمولی
منطق لئی
ایہ گوڈل
تے گینٹزن
نے اک دوجے
توں آزاد
دتا۔
فیر ساڈے
کول ایہ
اے:

\begin{thm}
\begin{enumerate}
\item $!A \liff !A^N$ کلاسیکی طور اُتے ثابت ہوندا اے۔
\item جے $!A$ کلاسیکی طور اُتے ثابت ہوئے، تاں
  $!A^N$ وجدانیت پسند طور اُتے ثابت ہوندا اے۔
\end{enumerate}
\end{thm}

ہن ایہ
گفتگو تصور
کرو۔ کلاسیکی
ریاضی دان:
”میں $!A$
ثابت کر
دتا!“
وجدانیت پسند
ریاضی دان:
”تہاڈا ثبوت
درست نہیں۔
اصل وچ تسی
$!A^N$ ثابت
کیتا اے۔“
کلاسیکی
ریاضی دان:
”مینوں ایہ
وی منظور!“
کلاسیکی
ریاضی دان
لئی وجدانیت
پسند بندہ
بس باریک
نکتے اُتے
اڑیا ہویا
اے، کیوں
جے اوہدے
حساب نال
دوویں ہم
ارز نیں۔
پر وجدانیت
پسند اصرار
کردا اے
کہ فرق
اے۔

چیتا رہے
کہ اوپرلا
ترجمہ صرف
خالص منطق
بارے اے؛
ایہ اس
سوال دا
جواب نہیں
دیندا کہ
کلاسیکی
تے وجدانیت
پسند ریاضی
لئی مناسب
\emph{غیر منطقی}
مسلّمے کیہڑے
نیں، یا
انہاں دا
باہمی رشتہ
کی اے۔
پر اوپرلے
قضیے دی
ہیٹھلی
چھوٹی
توسیع کجھ
فائدہ مند
خبر دیندی
اے:

\begin{thm}
جے $\Gamma$ کلاسیکی طور اُتے $!A$ ثابت کرے،
تاں $\Gamma^N$ وجدانیت پسند طور اُتے
$!A^N$ ثابت کردا اے۔
\end{thm}

دوجے لفظاں
وچ جے
$!A$ کجھ
مفروضیاں
توں کلاسیکی
طور اُتے
ثابت ہوئے،
تاں $!A^N$
انہاں مفروضیاں
دے دوہری
نفی والے
ترجمیاں
توں ثابت
ہوندا اے۔

اک جملے
یا قضیاتی
!!{formula}
دا وجدانیت
پسند اعتبار
وکھان لئی
ثبوت دینا
کافی اے۔
پر ایہ
کیویں وکھاؤ
گے کہ
اوہ معتبر
نہیں؟
ایہدے لئی
ایسی معنیات
چاہیدی اے
جیہڑی درست
ہوئے،
تے بہتر
اے تمام
وی ہوئے۔
کرپکے دی
معنیات ایس
لوڑ لئی
بہت ڈھکدی
اے۔

اسیں اوہی
کم کر
سکدے آں
جیہڑا کلاسیکی
منطق لئی
کیتا سی:
معنیات متعین
کرو، تے
درستگی تے
تمامیت ثابت
کرو۔
پر اک
فرق چیتے
رکھن جوگا
اے۔ کلاسیکی
منطق وچ
معنیات اک
لحاظ نال
”فطری“
سی، جیویں
منطقی رابطیاں
دے مطلب
وچ پہلے
توں موجود
ہوئے۔
کرپکے معنیات
لئی کجھ
بدیہی محرک
دتے جا
سکدے نیں،
پر اوہ
پہلی معنیات
وانگ ناگزیر
محسوس نہیں
ہوندی۔
ایس توں
علاوہ کلاسیکی
!!{structure}
اک فطری
ریاضیاتی تصور
اے؛ اسیں
یا تاں
!!a{structure}
نوں کلاسیکی
پہلے درجے
دی منطق
دے مطالعے
دا آلہ
سمجھ سکدے
آں، یا
کلاسیکی پہلے
درجے دی
منطق نوں
ریاضیاتی
!!{structure}s
دے مطالعے
دا آلہ۔
ایس دے
مقابلے وچ
کرپکے دیاں
!!{structure}s
صرف منطقی
بناوٹ ونجوں
دیکھیاں جا
سکدیاں
نیں؛ انہاں
دی کوئی
آزاد ریاضیاتی
دلچسپی نظر
نہیں آندی۔

قضیاتی زبان
لئی کرپکے
دی !!{structure}~
$\mModel M = \tuple{W, R, V}$
وچ اک
سیٹ~$W$،
جزوی ترتیب~$R$
جیہڑی~$W$
اُتے ہوئے
تے جیہدا
سب توں
چھوٹا !!{element}
ہوئے،
تے قضیاتی
متغیراں نوں
~$W$ دے
!!{element}s
نال ”سچ نوں برقرار رکھن والی“
(monotone) طور
تے جوڑن
والی قدر
بندی ہوندی
اے۔ خیال
ایہ اے
کہ $W$
دے !!{element}s
”جہاناں“
یا ”علمی
حالَتاں“ دی
نمائندگی کردے
نیں؛ اک
عنصر $v \geq
u$،~$u$
دی ”ممکنہ
اگلی حالت“
دسدا اے؛
تے~$u$
نال جوڑے
قضیاتی متغیر
اوہ قضیے
نیں جیہڑے
~$u$ دی
حالت وچ
سچے معلوم
نیں۔ ایس شرط دا مطلب اے کہ جیہڑا قضیاتی متغیر
اک حالت وچ سچا اے، اوہ اوس دی ہر ممکنہ اگلی
حالت وچ وی سچا رہندا اے؛ صرف متغیراں نوں
حالَتاں نال من چاہے طریقے نال جوڑنا کافی نہیں۔ فیر
تحمیلی صدق
(forcing)
دا تعلق
$\mSat{M}{!A}[w]$
ایس رشتے
نوں زبان
دے سارے
!!{formula}s
تک ودھاؤندا
اے؛
$\mSat{M}{!A}[w]$
نوں ”$!A$
حالت~$w$
وچ سچ
اے“ پڑھو۔
ایہ رشتہ
استقرائی
طور اُتے
ایویں متعین
اے:
\begin{enumerate}
\item $\mSat{M}{\Obj p_i}[w]$ جے تے صرف جے
  $\Obj p_i$ اوہناں قضیاتی متغیراں وچوں
  ہوئے جیہڑے~$w$ نال جوڑے گئے نیں۔
\item $\mSat/{M}{\lfalse}[w]$۔
\item $\mSat{M}{(!A \land !B)}[w]$ جے تے صرف جے
  $\mSat{M}{!A}[w]$ تے $\mSat{M}{!B}[w]$۔
\item $\mSat{M}{(!A \lor !B)}[w]$ جے تے صرف جے
  $\mSat{M}{!A}[w]$ یا $\mSat{M}{!B}[w]$۔
\item $\mSat{M}{(!A \lif !B)}[w]$ جے تے صرف جے
  جدوں وی $w' \geq w$ تے
  $\mSat{M}{!A}[w']$ ہون، تاں
  $\mSat{M}{!B}[w']$ وی ہوئے۔
\end{enumerate}
چنگی مشق
ایہ اے
کہ ایداں
دی کرپکے
!!{structure}
بنا کے،
جیہڑی مخالف
مثال دیوے،
وکھاؤ کہ
$\lnot (p \land q) \lif
(\lnot p \lor \lnot q)$
وجدانیت پسند
طور اُتے
معتبر نہیں۔

\end{document}
